\documentclass[11pt,fontset=fandol]{ctexart}
\usepackage[a4paper,margin=23mm]{geometry}
\usepackage{amsmath,amssymb,mathtools,bm,mathrsfs,esint}
\usepackage{graphicx,longtable,booktabs,array,calc}
\usepackage[normalem]{ulem}
\usepackage{xcolor,hyperref}
\hypersetup{hidelinks}
\setcounter{secnumdepth}{-1}
\newcounter{none}
\providecommand{\tightlist}{\setlength{\itemsep}{0pt}\setlength{\parskip}{0pt}}
\providecommand{\passthrough}[1]{#1}
\usepackage{amssymb,mathrsfs,esint,bm,newunicodechar}
\allowdisplaybreaks
\setlength{\emergencystretch}{3em}
\providecommand{\llangle}{\langle\!\langle}
\providecommand{\rrangle}{\rangle\!\rangle}
\newunicodechar{∞}{\ensuremath{\infty}}
\newunicodechar{∈}{\ensuremath{\in}}
\newunicodechar{∑}{\ensuremath{\sum}}
\newunicodechar{∫}{\ensuremath{\int}}
\newunicodechar{≤}{\ensuremath{\leq}}
\newunicodechar{≥}{\ensuremath{\geq}}
\newunicodechar{≠}{\ensuremath{\neq}}
\newunicodechar{→}{\ensuremath{\to}}
\newunicodechar{⇒}{\ensuremath{\Rightarrow}}
\newunicodechar{α}{\ensuremath{\alpha}}
\newunicodechar{β}{\ensuremath{\beta}}
\newunicodechar{γ}{\ensuremath{\gamma}}
\newunicodechar{λ}{\ensuremath{\lambda}}
\newunicodechar{φ}{\ensuremath{\varphi}}
\newunicodechar{ε}{\ensuremath{\varepsilon}}
\newunicodechar{σ}{\ensuremath{\sigma}}
\newunicodechar{θ}{\ensuremath{\theta}}
\newunicodechar{δ}{\ensuremath{\delta}}
\newunicodechar{η}{\ensuremath{\eta}}
\newunicodechar{π}{\ensuremath{\pi}}
\newunicodechar{ξ}{\ensuremath{\xi}}
\newunicodechar{Σ}{\ensuremath{\Sigma}}
\newunicodechar{Δ}{\ensuremath{\Delta}}
\newunicodechar{∂}{\ensuremath{\partial}}
\newunicodechar{∀}{\ensuremath{\forall}}
\newunicodechar{∃}{\ensuremath{\exists}}
\newunicodechar{√}{\ensuremath{\surd}}
\newunicodechar{−}{\ensuremath{-}}
\newunicodechar{μ}{\ensuremath{\mu}}
\newunicodechar{ρ}{\ensuremath{\rho}}
\newunicodechar{τ}{\ensuremath{\tau}}
\newunicodechar{ω}{\ensuremath{\omega}}
\newunicodechar{Ω}{\ensuremath{\Omega}}
\newunicodechar{χ}{\ensuremath{\chi}}
\newunicodechar{ψ}{\ensuremath{\psi}}
\newunicodechar{Γ}{\ensuremath{\Gamma}}
\newunicodechar{Φ}{\ensuremath{\Phi}}
\newunicodechar{Λ}{\ensuremath{\Lambda}}
\newunicodechar{ℝ}{\ensuremath{\mathbb{R}}}
\newunicodechar{ℤ}{\ensuremath{\mathbb{Z}}}
\newunicodechar{ℕ}{\ensuremath{\mathbb{N}}}
\newunicodechar{△}{\ensuremath{\triangle}}
\newunicodechar{⊗}{\ensuremath{\otimes}}
\begin{document}
\section{2019年杭州电子科技大学601数学分析考研真题}\label{ux5e74ux676dux5ddeux7535ux5b50ux79d1ux6280ux5927ux5b66601ux6570ux5b66ux5206ux6790ux8003ux7814ux771fux9898}

来源：\texttt{HDU/601数学分析/2019/2019年杭州电子科技大学601数学分析考研真题.pdf}。原
PDF 页序：1, 2。

\section{杭州电子科技大学}\label{ux676dux5ddeux7535ux5b50ux79d1ux6280ux5927ux5b66}

\section{2019年攻读硕士学位研究生招生考试}\label{ux5e74ux653bux8bfbux7855ux58ebux5b66ux4f4dux7814ux7a76ux751fux62dbux751fux8003ux8bd5}

\section{《数学分析》试题}\label{ux6570ux5b66ux5206ux6790ux8bd5ux9898}

（试题共4大题，共2页，总分150分）

姓名 报考专业 准考证号

【所有答案必须写在答题纸上，做在试卷或草稿纸上无效！】

\subsection{一、填空题（本大题共8小题，每小题5分，本大题共40分）}\label{ux4e00ux586bux7a7aux9898ux672cux5927ux9898ux51718ux5c0fux9898ux6bcfux5c0fux98985ux5206ux672cux5927ux9898ux517140ux5206}

1.极限
\(\lim _ { n \to \infty } \left( { \frac { 1 } { n + { \sqrt { 1 } } } } + { \frac { 1 } { n + { \sqrt { 2 } } } } + \cdots + { \frac { 1 } { n + { \sqrt { n } } } } \right) = (\)
）.

2.设g(0)=g'(0)=0,
\(f ( x ) = { \left\{ \begin{array} { l l } { g ( x ) \sin { \frac { 1 } { x } } , } & { x \neq 0 } \\ { 0 , } & { x = 0 } \end{array} \right. }\)
，则f(0)=（ ）.

3.曲线 \(L \left\{ { x = 1 - t ^ { 2 } } \right.\)
在t=1所对应点处的法线方程为（ ）.

4.曲线 \(y = x ^ { 4 } - 6 x ^ { 3 } + 1 2 x ^ { 2 } - x + 1\)
的拐点有（ ）.

5．设连续函数f(x)满足
\(2 f ( x ) + x = 1 + \int _ { 0 } ^ { 1 } f ( x ) d x\) ，则
\(\int _ { 0 } ^ { 1 } f ( x ) d x = (\) ）.

6.设
\(F ( y ) = \int _ { y } ^ { y ^ { 2 } } \frac { \sin ( x y ) } { x } d x .\)
则F(y)=（ ）.

7．设f(x,y)连续，则累次积分
\(\int _ { 1 } ^ { e ^ { 2 } } d y \int _ { \mathrm { \scriptsize ~ l n y } } ^ { 2 } f ( x , y ) d x\)
交换积分次序后的表达式为（ ）

8.若L为圆周
\(\left\{ { \begin{array} { l } { x ^ { 2 } + y ^ { 2 } + z ^ { 2 } = a ^ { 2 } } \\ { x + y + z = 0 } \end{array} } \right.\)
，则曲线积分
\(\int _ { L } ( 2 x ^ { 2 } + y ^ { 2 } + 2 x + y + 1 ) d s =\) （ ）：

\subsection{二、计算题（本大题共5小题，每小题6分，本大题共30分）}\label{ux4e8cux8ba1ux7b97ux9898ux672cux5927ux9898ux51715ux5c0fux9898ux6bcfux5c0fux98986ux5206ux672cux5927ux9898ux517130ux5206}

1．求极限值
\(\lim _ { x  0 ^ { + } } ( { \frac { \sin x } { x } } ) ^ { \frac { 1 } { x ^ { 2 } } }\)

2.设y=y(x)为 \(y = - y e ^ { x } + 2 e ^ { y } \sin x + 2 x\)
所确定的可导函数．求y(0).

3.设
\(f ( x ) = \left\{ { \begin{array} { l l } { x e ^ { - x ^ { 2 } } , } & { x \geq 0 } \\ { 0 , } & { x < 0 } \end{array} } \right.\)
，求 \(\int _ { 1 } ^ { x } f ( u - 1 ) d u\) 的表达式.

4.求级数
\(\sum _ { n = 1 } ^ { \infty } n \cdot \frac { 2 ^ { \frac { n } { 2 } } } { 3 ^ { n } }\)
的和.

5.计算 \(\iint _ { D } ( x + y + 1 ) ^ { 2 } d x d y ,\) 其中
\(D = \{ ( x , y ) | 1 \leq x ^ { 2 } + y ^ { 2 } \leq 4 \}\)

\subsection{三、综合题（本大题共3题，每题10分，本大题共30分)}\label{ux4e09ux7efcux5408ux9898ux672cux5927ux9898ux51713ux9898ux6bcfux989810ux5206ux672cux5927ux9898ux517130ux5206}

1.求a的值，使得
\(f ( x , y ) = \left\{ \begin{array} { l l } { { ( x ^ { 2 } + y ^ { 2 } ) ^ { a } \cos \frac { 1 } { \left| x \right| + \left| y \right| } , } } & { { ( x , y ) \neq ( 0 , 0 ) } } \\ { { 0 , } } & { { ( x , y ) = ( 0 , 0 ) } } \end{array} \right.\)
在点(0,0)可微.

2．讨论级数
\(\sum _ { n = 1 } ^ { \infty } { \frac { ( x + 1 ) ^ { n } } { ( n + 1 ) 2 ^ { n } } }\)
的敛散性（何时绝对收敛、何时条件收敛及何时发散）.

3．计算曲线积分
\(\int _ { L } { \frac { ( e ^ { x } - 2 x y ) d x - ( x ^ { 2 } + y ^ { 2 } ) d y } { ( x - 1 ) ^ { 2 } + y ^ { 2 } } }\)
，其中L是在圆周 \(y = \sqrt { 2 x - x ^ { 2 } }\)
上由点O(0,0)到点A(1,1)的一段弧.

\subsection{四、证明题（本大题共5题，每题10分，本大题共50分）}\label{ux56dbux8bc1ux660eux9898ux672cux5927ux9898ux51715ux9898ux6bcfux989810ux5206ux672cux5927ux9898ux517150ux5206}

1．设连续函数f(x)在{[}1,+∞)上为正值，且单调递减．令
\(a _ { n } = \sum _ { i = 1 } ^ { n } f ( i ) - \int _ { 1 } ^ { n } f ( x ) d x\)
证明 \(\{ a _ { n } \}\) 收敛.

2．设函数f(x)在区间{[}a,b{]}上二阶可导，曲线y=f(x)与过点
\(A ( a , f ( a ) )\) 和点\(B ( b , f ( b ) )\) 的直线相交于一点
\(C ( c , f ( c ) )\) ，其中 \(a < c < b\) ．试证：（1）至少存在一点
\(\xi _ { 1 } \in ( a , c )\) 及一点 \(\xi _ { 2 } \in ( c , b )\)
，使得
\(f ^ { \prime } ( { \boldsymbol { \xi } } _ { 1 } ) = f ^ { \prime } ( { \boldsymbol { \xi } } _ { 2 } )\)
；（2）至少存在一点 \(\xi \in ( a , b )\) ,使得
\(f ^ { \prime } ( \xi ) = 0\)

3．设1\textless a\textless b．证明函数
\(F ( x ) = \int _ { e } ^ { + \infty } { \frac { \cos t } { t ^ { x } } } d t\)
在{[}a,b{]}连续且可微.

4.设 \(\sum _ { i = 1 } ^ { \infty } a _ { n }\) 绝对收敛，证明
\(\sum _ { i = 1 } ^ { \infty } a _ { n } ( a _ { 1 } + a _ { 2 } + \cdots + a _ { n } )\)
也绝对收敛.

5．设R\textgreater0，Ω为曲面 \(z = R ^ { 2 } - x ^ { 2 } - y ^ { 2 }\)
与平面z=0所围成的空间闭区域，Σ为Ω表面的外侧，V为Ω的体积．证明
\(\iint x ^ { 2 } y e ^ { \sin z } d y d z - x y ^ { 2 } e ^ { \sin z } d x d z + z ( 1 + 2 x y z ) d x d y = V\)
\end{document}
